\centerline{\bf \Huge Вступительный разнобой}

\begin{flushright}
        \scriptsize{— Ты совсем не изменился.\\— Ну, конечно, я изменился. Ведь жить — значит меняться.\\Мисато и Кадзи}
\end{flushright}

\problem{1}
Найдите три нецелых положительных числа $a, b, c$ таких, что все числа

$$
\frac{a+b}{a-b}, \quad \frac{b+c}{b-c}, \quad \frac{c+a}{c-a}
$$

являются целыми.

\problem{2}
В начале года каждому из 150 бойцов лиги смешанных единоборств был присвоен номер от 1 до 150 . В течение года было проведено 149 поединков: первого со вторым, второго с третьим, ..., 149-го со 150 -м. В конце года был составлен список бойцов, победивших во всех поединках, в которых они участвовали в прошедшем году. Могли ли в этом списке оказаться и все бойцы с номерами кратными 17 , и все бойцы с номерами кратными $20 ?$

\problem{3}
Первоначально имеется один кусок сыра. Разрешается взять любой кусок сыра и проделать с ним одну из трех операций:

--- разделить его на два куска одинакового веса,

--- разделить его на 11 кусков одинакового веса,

--- разделить его на 23 куска одинакового веса.

Можно ли, используя только эти операции, разделить его на 2024 части одинакового веса?

\problem{4}
Найдутся ли такие 15 идущих подряд целых чисел, что сумма четырех наименьших из них равна сумме одиннадцати остальных?

\problem{5}
Можно ли число 240 представить в виде суммы девяти двузначных чисел (среди которых могут быть и одинаковые), в десятичной записи каждого из которых есть девятка?